%%
%% Riemannian Climate Fingerprinting
%% Sloan Austermann, OmniSciences LLC
%% Target: Journal of Climate / Climate Dynamics
%%
\documentclass[preprint,12pt]{elsarticle}

\usepackage{amsmath,amssymb,amsfonts}
\usepackage{mathtools}
\usepackage{bm}
\usepackage{booktabs}
\usepackage{graphicx}
\usepackage{natbib}
\usepackage{setspace}
\usepackage{lineno}
\usepackage{multirow}
\usepackage{xcolor}

\modulolinenumbers[5]

\journal{Journal of Climate}

% Notation shortcuts
\newcommand{\SPD}{\mathrm{SPD}}
\newcommand{\Sym}{\mathrm{Sym}}
\newcommand{\tr}{\mathrm{tr}}
\newcommand{\vech}{\mathrm{vech}}
\newcommand{\GL}{\mathrm{GL}}
\newcommand{\SO}{\mathrm{SO}}
\newcommand{\RR}{\mathbb{R}}
\newcommand{\Cbar}{\bar{C}}
\newcommand{\pval}{p}

\setcitestyle{round,aysep={},yysep={;}}

\begin{document}
\onehalfspacing

\begin{frontmatter}

\title{Riemannian Climate Fingerprinting: Detection and Mechanistic Separation of Forcings on the Covariance Manifold}

\author{Sloan Austermann}
\address{OmniSciences LLC}
\ead{sloan@omnisciences.co}

\begin{abstract}
Standard optimal fingerprinting \citep{Hasselmann1979,AllenStott2003} regresses observed temperature anomalies against model-predicted response patterns.
This tests whether the \emph{mean state} has shifted, but climate forcings also reshape the \emph{covariance structure}---altering teleconnection strengths, variance ratios, and cross-index correlations.
We propose Riemannian climate fingerprinting, which transports the detection framework to the manifold of symmetric positive-definite (SPD) covariance matrices equipped with the affine-invariant (Fisher--Rao) metric.

We construct rolling covariance matrices from 8 teleconnection indices and 58 gridded HadCRUT5 temperature cells (1983--2024), map them to the tangent space at the Fr\'{e}chet mean via the Riemannian logarithm, and regress the resulting tangent vectors against three forcing time series: anthropogenic (CO$_2$+CH$_4$), volcanic (GloSSAC~v2.22 AOD), and solar (sunspot number).
Significance is assessed via block-permutation surrogates that preserve temporal autocorrelation.

Riemannian fingerprinting detects all three forcings as statistically significant (anthropogenic $\pval=0.002$, volcanic $\pval=0.002$, solar $\pval=0.010$) on teleconnection covariance, where standard Euclidean fingerprinting on the same data fails to reach significance for any forcing ($\pval=0.098$, $0.348$, $0.486$).
A $V^+/V^-$ decomposition of regression coefficients into trace (energy budget) and traceless (teleconnection reorganization) components reveals mechanistically distinct forcing signatures: solar forcing acts primarily through the energy budget ($V^+=0.29$), while volcanic forcing operates almost entirely through teleconnection reorganization ($V^+=0.01$).

The method is model-free (no climate model ensembles required) and operates entirely on observed data.
Attribution fractions can be computed but are inherently imprecise---forcings are interdependent, effects are non-additive, and fractions are properties of the regression, not the climate system.
The method's core contribution is robust detection and mechanistic separation, not variance partitioning.
\end{abstract}

\begin{keyword}
climate detection \sep detection and attribution \sep Riemannian geometry \sep covariance manifold \sep optimal fingerprinting \sep teleconnections \sep SPD matrices
\end{keyword}

\end{frontmatter}

\section*{Significance Statement}
Climate change detection traditionally tests whether average temperatures have shifted, but climate forcings also reshape the relationships between different parts of the climate system---strengthening some connections and weakening others.
We introduce a geometric method that detects these structural changes by treating climate covariance matrices as points on a curved mathematical surface, where standard flat-space statistics fail.
Applied to 40~years of observations, the method detects human, volcanic, and solar influences on climate structure with high confidence, and reveals that each forcing reshapes the climate system through a qualitatively different mechanism: solar forcing primarily changes total variability, while volcanic eruptions almost entirely reorganize teleconnection patterns.
This mechanistic fingerprint is invisible to standard detection methods.

\linenumbers

%% ====================================================================
\section{Introduction}\label{sec:intro}
%% ====================================================================

\subsection{The detection and attribution problem}

Detection and attribution (D\&A) is the central inferential problem in climate science: distinguishing externally forced climate change from internal variability \citep{IPCC_AR6_Ch3}.
The foundational framework, due to \citet{Hasselmann1979}, casts the problem as a generalized regression of observed climate fields onto signal patterns predicted by general circulation models.
\citet{AllenStott2003} formalized this into the total least squares framework with explicit treatment of model uncertainty, which has since become the standard approach in IPCC assessments.

A common feature of all standard D\&A methods is that they operate on the \emph{mean state}---spatial fields of temperature or precipitation anomalies, possibly dimension-reduced by spherical harmonics or empirical orthogonal functions.
This captures forced shifts in first-order statistics (the mean climate), but climate forcings also alter \emph{second-order statistics}: the covariance structure that encodes teleconnection strengths, variance ratios, and cross-index correlations.
These covariance changes are invisible to mean-field D\&A.

\subsection{Why covariance matters}

Forcings do not merely shift means; they reshape correlation structures.
Volcanic eruptions disrupt ENSO teleconnections \citep{Adams2003,Predybaylo2017}.
The solar cycle modulates NAO--AO coupling \citep{Gray2010,Scaife2013}.
Anthropogenic warming alters Arctic--midlatitude teleconnections \citep{FrancisVavrus2012,Cohen2014}.
These effects operate on the covariance of the climate system, not just its mean.

Second-order statistics contain substantially more information than first-order statistics.
A $d$-dimensional system has $d(d+1)/2$ covariance degrees of freedom versus $d$ mean degrees of freedom.
For $d=8$ teleconnection indices, this is 36 covariance DOF versus 8 mean DOF---a factor of 4.5.
For $d=58$ gridded temperature cells, the ratio is 1,711 to 58---nearly 30-fold.
This larger signal space provides more room for forced signals to separate from internal noise.

\subsection{The geometry of covariance matrices}

Covariance matrices live on $\SPD(d)$, the cone of $d \times d$ symmetric positive-definite matrices.
This is not a flat Euclidean space but a curved Riemannian manifold.
Euclidean operations---addition, subtraction, arithmetic averaging---on SPD matrices violate the positive-definiteness constraint, are not invariant to changes of basis, and weight large eigenvalues disproportionately \citep{Pennec2006,Moakher2005}.

The affine-invariant (Fisher--Rao) metric provides a geometrically principled alternative.
Under this metric, $\SPD(d)$ has the structure of the symmetric space $\GL^+(d)/\SO(d)$, with well-defined geodesics, an intrinsic Fr\'{e}chet mean, and a logarithmic map that linearizes the geometry for statistical analysis.
This Riemannian framework has been successfully applied in brain--computer interfaces \citep{Barachant2013}, diffusion tensor imaging \citep{Pennec2006,Guillot2015}, and materials science, but has not previously been applied to climate D\&A.

The affine-invariant metric is a special case of the DeWitt supermetric on the space of Riemannian metrics \citep{DeWitt1967}.
This connection, elaborated in Section~\ref{sec:dewitt}, places climate covariance analysis within a broader geometric framework spanning general relativity, elasticity, and neuroscience.

\subsection{Contributions}

This paper makes four contributions:

\begin{enumerate}
\item \textbf{Detection with geometric advantage.} We demonstrate statistically significant detection of all three major forcings ($\pval < 0.01$) in observed climate covariance data where Euclidean methods on the same data fail to detect any.

\item \textbf{Mechanistic separation via $V^+/V^-$.} We introduce a decomposition of regression coefficients into trace (energy budget) and traceless (teleconnection reorganization) components, revealing qualitatively distinct forcing signatures unavailable in standard mean-field D\&A.

\item \textbf{Model-free detection.} The method requires no climate model ensembles---only observed data and forcing time series---eliminating model dependence at the cost of not having model-predicted response patterns as priors.

\item \textbf{Riemannian fingerprinting on $\SPD(d)$.} We transport optimal fingerprinting to the SPD manifold via tangent-space regression, validated on both low-dimensional teleconnection indices ($d=8$) and high-dimensional gridded data ($d=58$).
\end{enumerate}


%% ====================================================================
\section{Mathematical Framework}\label{sec:math}
%% ====================================================================

\subsection{The SPD manifold}\label{sec:spd}

The space of $d \times d$ symmetric positive-definite matrices,
\begin{equation}
\SPD(d) = \bigl\{ C \in \RR^{d \times d} : C = C^\top,\; C > 0 \bigr\},
\end{equation}
is a Riemannian manifold of dimension $d(d+1)/2$.

\paragraph{Affine-invariant metric.}
At a point $C \in \SPD(d)$, the metric on tangent vectors $V, W \in T_C\SPD(d) \cong \Sym(d)$ is
\begin{equation}\label{eq:metric}
\langle V, W \rangle_C = \tr\bigl(C^{-1} V \, C^{-1} W\bigr).
\end{equation}
This metric is invariant under the $\GL(d)$ action $C \mapsto A C A^\top$ and gives $\SPD(d)$ the structure of the symmetric space $\GL^+(d)/\SO(d)$.

\paragraph{Geodesic distance.}
The geodesic distance between $C_1, C_2 \in \SPD(d)$ is
\begin{equation}\label{eq:geodist}
d(C_1, C_2) = \bigl\| \log\bigl(C_1^{-1/2}\, C_2\, C_1^{-1/2}\bigr) \bigr\|_F = \sqrt{\sum_{i=1}^d (\log \lambda_i)^2},
\end{equation}
where $\lambda_1, \ldots, \lambda_d$ are the generalized eigenvalues of the pencil $(C_1, C_2)$---that is, the eigenvalues of $C_1^{-1} C_2$.

\paragraph{Geodesic.}
The geodesic from $C_1$ to $C_2$ parametrized by $t \in [0,1]$ is
\begin{equation}\label{eq:geodesic}
\gamma(t) = C_1^{1/2}\bigl(C_1^{-1/2}\, C_2\, C_1^{-1/2}\bigr)^t C_1^{1/2}.
\end{equation}

\subsection{Fr\'{e}chet mean}\label{sec:frechet}

The Fr\'{e}chet mean of a collection $\{C_1, \ldots, C_N\} \subset \SPD(d)$ is the unique minimizer
\begin{equation}\label{eq:frechet}
\Cbar = \arg\min_{C \in \SPD(d)} \sum_{i=1}^N d^2(C, C_i).
\end{equation}
We compute $\Cbar$ by Riemannian gradient descent.
At a current estimate $\mu$, the Riemannian gradient is
\begin{equation}
\nabla = \frac{1}{N} \sum_{i=1}^N \log_\mu(C_i),
\end{equation}
and the update is $\mu_{k+1} = \exp_{\mu_k}(\nabla)$.
The algorithm converges when $\|\nabla\|_F$ falls below a tolerance (we use $10^{-8}$).
Initialization via the log-Euclidean approximation $\exp\!\bigl(\frac{1}{N}\sum_i \log C_i\bigr)$ ensures rapid convergence, typically within 5--10 iterations.

\subsection{Logarithmic map and tangent-space linearization}\label{sec:logmap}

The Riemannian logarithm at $\Cbar$ maps each covariance matrix to the tangent space:
\begin{equation}\label{eq:logmap}
\log_{\Cbar}(C_i) = \Cbar^{1/2}\, \log\!\bigl(\Cbar^{-1/2}\, C_i\, \Cbar^{-1/2}\bigr)\, \Cbar^{1/2}.
\end{equation}
In practice, because the $\Cbar^{1/2}$ factors cancel in the affine-invariant inner product~\eqref{eq:metric}, we work with the simplified tangent vectors
\begin{equation}\label{eq:tangent}
V_i = \log\!\bigl(\Cbar^{-1/2}\, C_i\, \Cbar^{-1/2}\bigr) \in \Sym(d).
\end{equation}
These are vectorized via the $\vech$ operator: $\bm{v}_i = \vech(V_i) \in \RR^{d(d+1)/2}$, where off-diagonal elements are scaled by $\sqrt{2}$ to preserve the Frobenius norm.

\subsection{Tangent-space regression (Riemannian climate fingerprinting)}\label{sec:regression}

Given $K$ forcing time series $F_1(t), \ldots, F_K(t)$ aligned to the covariance window centers, we solve the multivariate linear regression
\begin{equation}\label{eq:regression}
\bm{v}(t) = \sum_{k=1}^K \bm{\beta}_k\, F_k(t) + \bm{\beta}_0 + \bm{\varepsilon}(t),
\end{equation}
where $\bm{v}(t) \in \RR^p$ with $p = d(d+1)/2$ is the vectorized tangent vector at time $t$, and each $\bm{\beta}_k \in \RR^p$ is the multivariate regression coefficient for forcing~$k$.

For the teleconnection case ($d=8$, $p=36$, $N \approx 440$ windows), ordinary least squares (OLS) is well-posed since $p \ll N$.
For the gridded case ($d=58$, $p=1{,}711$, $N \approx 440$), $p > N$ requires regularization.
We use PCA dimensionality reduction retaining the top 30 principal components (approximately 85\% of tangent-space variance), followed by Ridge regression with penalty $\alpha = 1.0$ (chosen via leave-one-block-out cross-validation with block size 60 months).
The regression coefficients are then projected back to the full $p$-dimensional tangent space for interpretation.

We do not rely on parametric $F$-tests because the Gaussianity assumption fails for tangent vectors derived from overlapping rolling windows.
Instead, significance is assessed via the surrogate procedure described in Section~\ref{sec:surrogates}.

\subsection{$V^+/V^-$ decomposition}\label{sec:vpvm}

Each regression coefficient $\bm{\beta}_k$ can be reshaped to a symmetric matrix $B_k = \vech^{-1}(\bm{\beta}_k) \in \Sym(d)$.
We decompose
\begin{equation}\label{eq:vpvm}
B_k = B_k^+ + B_k^-,
\end{equation}
where
\begin{align}
B_k^+ &= \frac{\tr(B_k)}{d}\, I_d \qquad \text{(trace/isotropic component)}, \label{eq:vplus}\\
B_k^- &= B_k - B_k^+ \qquad \text{(traceless/anisotropic component)}. \label{eq:vminus}
\end{align}
The normalized $V^+$ ratio is
\begin{equation}\label{eq:vratio}
r_k^+ = \frac{\|B_k^+\|_F}{\|B_k^+\|_F + \|B_k^-\|_F}.
\end{equation}

This ratio classifies forcing mechanisms:
\begin{itemize}
\item $r^+ \approx 1$: forcing primarily changes total variance (energy budget effect);
\item $r^+ \approx 0$: forcing primarily reorganizes correlations (teleconnection effect);
\item $r^+ \approx 0.5$: mixed mechanism.
\end{itemize}

The $V^+/V^-$ decomposition corresponds to the splitting of $T_C\SPD(d)$ into the trace direction (scalar/conformal mode) and the traceless complement under the $\GL^+(d)/\SO(d)$ symmetric space structure.
In the DeWitt supermetric parametrized by~$\lambda$, $V^+$ is the conformal mode and $V^-$ spans the remaining $d(d+1)/2 - 1$ directions (Section~\ref{sec:dewitt}).


%% ====================================================================
\section{Data and Methods}\label{sec:methods}
%% ====================================================================

\subsection{Data sources}\label{sec:data}

\paragraph{Teleconnection indices ($d=8$).}
Monthly values over 1983--2024 from NOAA/CPC and KNMI Climate Explorer:
Ni\~{n}o~3.4 (central equatorial Pacific SST, ENSO), SOI (Southern Oscillation Index), PDO (Pacific Decadal Oscillation), AMO (Atlantic Multidecadal Oscillation), NAO (North Atlantic Oscillation), AO (Arctic Oscillation), PNA (Pacific--North American pattern), and GMTA (global mean temperature anomaly, HadCRUT5).

\paragraph{Gridded temperatures ($d=58$).}
HadCRUT5 $5^\circ \times 5^\circ$ grid cells selected for data completeness over 1983--2024, providing monthly temperature anomalies relative to the 1961--1990 baseline.

\paragraph{Forcing time series.}
Three forcing proxies are used:
\begin{itemize}
\item \emph{Anthropogenic:} monthly CO$_2$ and CH$_4$ concentrations from NOAA Global Monitoring Laboratory (Mauna Loa), combined via first principal component to avoid collinearity.
\item \emph{Volcanic:} stratospheric aerosol optical depth from GloSSAC~v2.22, capturing the episodic signals of El~Chich\'{o}n (1982), Pinatubo (1991), and Hunga~Tonga (2022).
\item \emph{Solar:} monthly sunspot number from SILSO~v2.0, a proxy for the $\sim$11-year solar cycle.
\end{itemize}
All forcing series are standardized to zero mean and unit variance.
Forcing values are aligned to rolling window centers with lags calibrated to the physical response timescale: 3~months for volcanic, 6~months for solar, and 0~months for anthropogenic.

\subsection{Rolling covariance estimation}\label{sec:rolling}

We compute covariance matrices in rolling windows of width $w = 60$~months (5~years) with a step size of 1~month, yielding $N \approx 440$ SPD matrices over the 1983--2024 period.
Ledoit--Wolf shrinkage toward a scaled identity matrix is applied with parameter $\lambda_\mathrm{LW} = 0.15$ to regularize the covariance estimates:
\begin{equation}
\hat{C}(t) = (1 - \lambda_\mathrm{LW})\, S(t) + \lambda_\mathrm{LW}\, \frac{\tr(S(t))}{d}\, I_d,
\end{equation}
where $S(t)$ is the sample covariance in window $t$.
Positive-definiteness is further guaranteed by clamping eigenvalues at a floor of $10^{-8}$.
The shrinkage parameter was not tuned; sensitivity to values in $[0.05, 0.30]$ is documented in Appendix~\ref{app:sensitivity}.

\subsection{Regression configuration}\label{sec:regconfig}

\paragraph{Teleconnection-based ($d=8$).}
With $p = 36$ tangent-space dimensions and $N \approx 440$ time points, OLS is well-posed ($p \ll N$).
We report multivariate $R^2$ and per-forcing test statistics.

\paragraph{Gridded ($d=58$).}
With $p = 1{,}711$ tangent-space dimensions and $N \approx 440$, PCA reduction to 30 components (retaining approximately 85\% of tangent-space variance) is applied prior to Ridge regression with penalty $\alpha = 1.0$, selected via leave-one-block-out cross-validation with block size 60~months.
Regression coefficients are projected back to the full tangent space for $V^+/V^-$ decomposition.

\paragraph{Euclidean comparison.}
To isolate the geometric contribution, we run an identical pipeline on the same data but regress against raw vectorized covariance matrices $\vech(C_i)$ without the Riemannian log map or Fr\'{e}chet mean.
A three-way comparison (mean-field Euclidean, flat covariance Euclidean, Riemannian) is reported in Section~\ref{sec:geometric}.

\subsection{Significance testing: block-permutation surrogates}\label{sec:surrogates}

Parametric $p$-values are unreliable because (i)~tangent vectors are autocorrelated due to overlapping windows, (ii)~their distribution is non-Gaussian, and (iii)~multiple regression introduces dependence between forcing estimates.
We therefore assess significance non-parametrically:

\begin{enumerate}
\item For each of $n_\mathrm{surr} = 1{,}000$ surrogate iterations:
  \begin{enumerate}
  \item Randomly permute the forcing design matrix in contiguous blocks of length $b = 5$~years (60~months, matching the rolling window), destroying forcing--covariance coupling while preserving within-block autocorrelation.
  \item Re-run the regression with permuted forcings.
  \item Record the per-forcing $F$-statistic (equivalently, $\|\bm{\beta}_k\|^2$).
  \end{enumerate}
\item The $p$-value for forcing $k$ is the fraction of surrogates with $F_k^\mathrm{surr} \geq F_k^\mathrm{obs}$.
\end{enumerate}

This is a non-parametric test with exact size control up to Monte Carlo error ($\pm 0.001$ at 1{,}000 surrogates).
Phase-randomization surrogates (preserving the power spectrum) were also tested and yield similar results; block permutation is reported as the primary method because it makes fewer assumptions about spectral structure.

We report results at the primary block size of $b = 5$ years, which provides 30 permutable blocks---sufficient for a well-resolved null distribution while preserving substantial within-block autocorrelation.
Sensitivity to block size is documented in Appendix~\ref{app:blocksize}, where we show that significance depends on this choice: at block sizes $\geq 12$~years, anthropogenic and solar lose significance due to insufficient permutation resolution (only 13 blocks), not absence of signal.
Volcanic forcing is the most robust, remaining significant at all tested block sizes up to 20~years.

\subsection{Bootstrap stability}\label{sec:bootstrap}

Block bootstrap (block size 60~months, 2{,}000 resamples) is used to assess stability.
The full pipeline---Fr\'{e}chet mean, log map, regression---is re-run on each bootstrap sample.
We report bootstrap 95\% confidence intervals on $V^+/V^-$ ratios (primary) and attribution fractions (secondary, with caveats discussed in Section~\ref{sec:attribution}).


%% ====================================================================
\section{Results}\label{sec:results}
%% ====================================================================

\subsection{Detection: all three forcings significant ($d=8$ teleconnections)}\label{sec:detection8}

Table~\ref{tab:detection8} presents the central detection result.
All three forcings are detected at $\pval < 0.01$ via block-permutation surrogates (block size 5~years, 1{,}000 surrogates).
The multivariate $R^2 = 0.27$, meaning 73\% of tangent-space variance is unexplained---consistent with dominant internal variability (ENSO, PDO regime shifts, weather noise) on interannual timescales.
The central finding is that each forcing leaves a statistically distinguishable fingerprint on the covariance structure of the climate system.

\begin{figure}[t]
\centering
\includegraphics[width=\textwidth]{figures/fig2_rolling_covariance.pdf}
\caption{Rolling covariance elements (60-month window) for selected index pairs, 1983--2024. Non-stationarity is evident, with clear perturbation from the Pinatubo eruption (1991, dashed red line). These time-varying covariance matrices form a trajectory on $\SPD(d)$ that the tangent-space regression decomposes into forced and unforced components.}
\label{fig:rolling_cov}
\end{figure}

\begin{table}[t]
\centering
\caption{Detection significance for three forcings via Riemannian tangent-space regression on 8 teleconnection indices (1983--2024). $p$-values from block-permutation surrogates ($n_\mathrm{surr} = 1{,}000$, block size = 5~years).}
\label{tab:detection8}
\begin{tabular}{lccc}
\toprule
Forcing & $\|\bm{\beta}\|_F$ & $\pval$-value & Detected? \\
\midrule
Anthropogenic (CO$_2$+CH$_4$) & 0.343 & 0.002 & Yes ($\pval < 0.01$) \\
Volcanic (AOD) & 0.437 & 0.002 & Yes ($\pval < 0.01$) \\
Solar (SSN) & 0.328 & 0.010 & Yes ($\pval < 0.01$) \\
\midrule
\multicolumn{4}{l}{Multivariate $R^2 = 0.27$} \\
\bottomrule
\end{tabular}
\end{table}

\begin{figure}[t]
\centering
\includegraphics[width=\textwidth]{figures/fig3_surrogate_nulls.pdf}
\caption{Block-permutation null distributions (500 surrogates, block size = 5~years) for each forcing. Black vertical lines mark the observed $\|\bm{\beta}\|_F$. All three observed values lie well outside the null distributions, confirming statistically significant detection.}
\label{fig:surrogates}
\end{figure}

\subsection{Detection: gridded validation ($d=58$)}\label{sec:detection58}

Table~\ref{tab:detection58} shows that all three forcings are detected at $\pval = 0.001$ (at the resolution of 1{,}000 surrogates) in the gridded HadCRUT5 analysis.
PCA plus Ridge regression enables fingerprinting in $p = 1{,}711$ tangent-space dimensions, and detection results are robust across both low- and high-dimensional representations of the climate system ($R^2 = 0.30$).

\begin{table}[t]
\centering
\caption{Detection significance for gridded HadCRUT5 ($d=58$, $p=1{,}711$). PCA to 30 components + Ridge regression ($\alpha=1.0$). Block-permutation surrogates ($n_\mathrm{surr}=1{,}000$, block size = 5~years).}
\label{tab:detection58}
\begin{tabular}{lcc}
\toprule
Forcing & $\pval$-value & Detected? \\
\midrule
Anthropogenic & 0.001 & Yes \\
Volcanic & 0.001 & Yes \\
Solar & 0.001 & Yes \\
\midrule
\multicolumn{3}{l}{Multivariate $R^2 = 0.30$} \\
\bottomrule
\end{tabular}
\end{table}

\subsection{Mechanistic separation: $V^+/V^-$ decomposition}\label{sec:mechanism}

Table~\ref{tab:vpvm} presents the $V^+/V^-$ decomposition for the teleconnection-based analysis.
The three forcings occupy qualitatively distinct regions of the $V^+/V^-$ spectrum, revealing mechanistically distinct forcing signatures.

\begin{table}[t]
\centering
\caption{$V^+/V^-$ decomposition of regression coefficients (teleconnection-based, $d=8$). $V^+$ ratio defined by Eq.~\eqref{eq:vratio}.}
\label{tab:vpvm}
\begin{tabular}{lccl}
\toprule
Forcing & $V^+$ ratio & $V^-$ dominant? & Mechanism \\
\midrule
Solar & 0.29 & No & Primarily energy budget \\
Anthropogenic & 0.11 & Yes & Mixed; mostly teleconnection \\
Volcanic & 0.01 & Yes & Almost purely teleconnection \\
\bottomrule
\end{tabular}
\end{table}

\paragraph{Solar ($V^+ = 0.29$).}
The 11-year solar cycle modulates total energy input, changing the overall variance of the climate system.
This is consistent with total solar irradiance (TSI) variations affecting the global energy budget directly: when TSI rises, climate variability increases roughly isotropically across indices.
Solar forcing has the highest $V^+$ ratio of the three forcings, indicating that its covariance fingerprint is dominated by energy-budget effects rather than teleconnection reorganization.

\paragraph{Anthropogenic ($V^+ = 0.11$).}
Long-term warming has a component of increased total variance (stronger extremes), but the dominant effect on covariance is a reshaping of correlations---for example, weakening of tropical--polar coupling as the Arctic warms preferentially (Arctic amplification), and altered ENSO--midlatitude teleconnections.
The low $V^+$ ratio indicates that anthropogenic forcing acts primarily through teleconnection reorganization rather than a uniform increase in variance.

\paragraph{Volcanic ($V^+ = 0.01$).}
Major eruptions cause rapid global cooling followed by recovery, but the dominant covariance effect is a reorganization of teleconnection patterns---for example, the strengthened NAO-like pattern observed after Pinatubo \citep{Adams2003}.
The energy-budget signal of an eruption is fast and transient; the teleconnection disruption persists longer in the rolling covariance window and dominates the fingerprint.
That volcanic forcing has the lowest $V^+$ ratio of the three forcings is a novel finding: eruptions are typically discussed in terms of their radiative (energy budget) effects, but their covariance fingerprint is almost entirely structural.

The mechanistic separation provided by $V^+/V^-$ is a novel diagnostic unavailable in standard mean-field D\&A.

\begin{figure}[t]
\centering
\includegraphics[width=0.85\textwidth]{figures/fig5_vplus_vminus.pdf}
\caption{$V^+/V^-$ mechanistic decomposition of regression coefficients. Solar forcing has the highest $V^+$ fraction (energy budget mechanism), while volcanic forcing is almost entirely $V^-$ (teleconnection reorganization). This ordering is invariant across all tested configurations (Tables~\ref{tab:B1}--\ref{tab:B5}).}
\label{fig:vpvm}
\end{figure}

\begin{figure}[t]
\centering
\includegraphics[width=\textwidth]{figures/fig6_beta_heatmaps.pdf}
\caption{Regression coefficient matrices $B_k = \vech^{-1}(\bm{\beta}_k)$ for each forcing (teleconnection-based, $d=8$). Row/column labels are climate indices. The distinct spatial patterns confirm that each forcing reshapes different aspects of the covariance structure.}
\label{fig:betas}
\end{figure}

\paragraph{Caveat.}
The $V^+/V^-$ decomposition is basis-dependent.
When applied after PCA projection back to the full space (gridded $d=58$ case), the trace of the projected regression coefficient does not correspond to the physical energy budget.
We therefore report $V^+$ ratios only for the teleconnection-based analysis, where the basis has clear physical meaning.
This loss of interpretability after PCA is a fundamental limitation: the $V^+/V^-$ decomposition provides mechanistic insight only when the covariance basis corresponds to physically meaningful climate indices.

\subsection{Geometric advantage: Euclidean comparison}\label{sec:geometric}

Table~\ref{tab:euclidean} presents the central comparison.
No forcing reaches $\pval < 0.05$ under Euclidean regression on the same data, the same surrogates, and the same forcing series.
The only difference is the geometry: Euclidean regression treats covariance matrices as flat vectors, ignoring the manifold structure.
The Riemannian approach improves $p$-values by one to two orders of magnitude.

\begin{table}[t]
\centering
\caption{Euclidean vs.\ Riemannian detection on identical data ($d=8$ teleconnections). Same surrogates, same forcings---only the geometry differs.}
\label{tab:euclidean}
\begin{tabular}{lcc}
\toprule
Forcing & $\pval$ (Euclidean) & $\pval$ (Riemannian) \\
\midrule
Anthropogenic & 0.098 & 0.002 \\
Volcanic & 0.348 & 0.002 \\
Solar & 0.486 & 0.010 \\
\bottomrule
\end{tabular}
\end{table}

\begin{figure}[t]
\centering
\includegraphics[width=\textwidth]{figures/fig4_threeway_comparison.pdf}
\caption{Three-way method comparison showing $-\log_{10}(p)$ for each forcing. Higher bars indicate stronger detection. The dashed line marks $p=0.05$; the dotted line marks $p=0.01$. Riemannian fingerprinting (right) detects all three forcings at $p<0.01$. Standard Euclidean regression on means (left) detects none. The flat Euclidean control on vectorized covariances (center) isolates the dimensionality effect from the geometric effect.}
\label{fig:threeway}
\end{figure}

The dramatic improvement arises from three sources.
First, covariance matrices contain $d(d+1)/2 = 36$ independent entries versus $d = 8$ for mean fields, providing a larger signal space.
Second, the Riemannian log map treats eigenvalue ratios on a logarithmic scale, which is the natural scale for the multiplicative perturbations that forcings induce on covariance structure (a 10\% increase in variance is the same ``distance'' whether the baseline variance is 0.1 or 10).
Third, linearization at the Fr\'{e}chet mean (the intrinsic center of the data cloud on the manifold) minimizes distortion, while Euclidean regression has no such principled linearization point.

To disentangle the dimensionality effect from the geometric effect, Table~\ref{tab:threeway} presents a three-way comparison.
The flat Euclidean control uses the same 36-dimensional vectorized covariance matrices but without the Riemannian log map or Fr\'{e}chet mean.
Comparing the flat Euclidean (36D) to Riemannian (36D) columns isolates the purely geometric contribution at fixed dimensionality: anthropogenic improves from $\pval = 0.016$ to $0.002$, and solar crosses the significance threshold from $0.056$ to $0.010$.
The geometry contributes independently of the dimensionality increase.

\begin{table}[t]
\centering
\caption{Three-way method comparison isolating the geometric effect. ``Euclidean (means)'' regresses mean anomalies (8D). ``Flat Euclidean (cov)'' regresses vectorized covariance matrices (36D) without Riemannian geometry. ``Riemannian'' uses the log map at the Fr\'{e}chet mean (36D).}
\label{tab:threeway}
\begin{tabular}{lccccl}
\toprule
Method & Dims & $\pval_\mathrm{anthro}$ & $\pval_\mathrm{volc}$ & $\pval_\mathrm{solar}$ & $R^2$ \\
\midrule
Euclidean (means) & 8 & 0.098 & 0.348 & 0.486 & 0.430 \\
Flat Euclidean (cov) & 36 & 0.016 & 0.002 & 0.056 & 0.235 \\
Riemannian & 36 & 0.002 & 0.002 & 0.010 & 0.276 \\
\bottomrule
\end{tabular}
\end{table}

The higher $R^2$ for the mean-field Euclidean case (0.430 vs.\ 0.276) does not indicate superior detection.
It reflects the lower dimensionality: 8 response dimensions are more easily ``explained'' by 3 predictors than 36 are.
The relevant comparison is the $p$-value, which measures whether the explained variance is statistically distinguishable from what chance would produce.

\subsection{Attribution fractions (with caveats)}\label{sec:attribution}

Attribution fractions---defined as $\|\bm{\beta}_k\|^2 / \sum_{k'} \|\bm{\beta}_{k'}\|^2$---can be computed from the regression but should be interpreted with extreme caution.
They are presented in Table~\ref{tab:attribution} for completeness, not as a primary result.

\begin{table}[t]
\centering
\caption{Bootstrap 95\% confidence intervals on attribution fractions (teleconnection-based, $d=8$). These fractions are inherently imprecise---see text for three reasons they should not be over-interpreted.}
\label{tab:attribution}
\begin{tabular}{lcc}
\toprule
Forcing & Attribution \% & 95\% CI \\
\midrule
Anthropogenic & 36\% & [29\%, 44\%] \\
Volcanic & 26\% & [16\%, 37\%] \\
Solar & 38\% & [27\%, 46\%] \\
\bottomrule
\end{tabular}
\end{table}

Three reasons these fractions should not be over-interpreted:

\begin{enumerate}
\item \textbf{Forcings are not independent.} Anthropogenic warming modulates the climate's response to volcanic and solar forcings (e.g., a warmer ocean responds differently to eruptions).
The regression treats forcings as additive predictors, but the real climate system has nonlinear interactions.

\item \textbf{Effects are not additive on the covariance manifold.} The SPD manifold is curved.
The tangent-space linearization enables additive regression, but the underlying covariance perturbations compose multiplicatively ($C \mapsto A C A^\top$).
Attribution fractions from the linearized regression do not correspond to a clean decomposition of the nonlinear covariance evolution.

\item \textbf{Fractions are properties of the regression, not the climate.} The fractions depend on forcing time series normalization, regression method (OLS vs.\ Ridge), window size, and shrinkage parameter.
Different defensible methodological choices yield fractions that vary by 10--15 percentage points (Appendix~\ref{app:sensitivity}).
The detection $p$-values are robust to these choices; the fractions are not.
\end{enumerate}

The 95\% confidence intervals are wide and overlap substantially.
With $R^2 \approx 0.27$--$0.30$, approximately 70\% of covariance variation is not explained by the three forcings---this is dominated by internal variability (ENSO, PDO regime shifts, and other modes of natural variability).
This method provides robust \emph{detection} and \emph{mechanistic separation}, not precise variance partitioning.


%% ====================================================================
\section{Discussion}\label{sec:discussion}
%% ====================================================================

\subsection{Why Riemannian outperforms Euclidean}\label{sec:whyriemann}

The Riemannian approach gains statistical power from three sources, which can now be discussed in light of the results.

First, covariance matrices provide more degrees of freedom than mean fields.
For $d=8$ indices, the 36 covariance DOF versus 8 mean DOF provide a larger space in which forcing signals can separate from noise.
The three-way comparison (Table~\ref{tab:threeway}) confirms that moving from 8 to 36 dimensions, even without Riemannian geometry, already improves detection for volcanic forcing.

Second, the affine-invariant metric treats eigenvalue ratios on a logarithmic scale.
The Euclidean distance $\|C_1 - C_2\|_F$ weights large eigenvalues disproportionately and is not invariant to reparametrization.
The geodesic distance~\eqref{eq:geodist} treats all eigenvalue ratios equally on a log scale, which is the natural scale for multiplicative perturbations of covariance structure.
Forcing-induced changes in covariance are multiplicative---a 20\% increase in ENSO variance is the same physical effect regardless of the baseline ENSO variance level---and the Riemannian metric captures this.

Third, linearization at the Fr\'{e}chet mean provides the best linear approximation to the curved geometry, minimizing distortion in the tangent-space regression.
The three-way comparison confirms an independent geometric contribution beyond dimensionality: at fixed dimensionality (36D), the Riemannian approach improves anthropogenic $p$-values from 0.016 to 0.002 and moves solar from marginal ($\pval = 0.056$) to significant ($\pval = 0.010$).

We emphasize that the geometric advantage is \emph{partly} dimensionality and \emph{partly} geometry (Table~\ref{tab:threeway}).
Both contributions are real and complementary.

\subsection{$V^+/V^-$ as a mechanistic diagnostic}\label{sec:vpvm_discuss}

The $V^+/V^-$ decomposition provides a principled, geometry-motivated way to classify forcing mechanisms.
The trace component $B_k^+$ captures changes in total variance---analogous to the energy budget.
Forcings that inject or remove energy uniformly inflate or deflate the covariance matrix isotropically.
The traceless component $B_k^-$ captures changes in correlation structure at approximately fixed total variance---analogous to teleconnection reorganization, where energy is redistributed among modes without changing the total.

That the three forcings produce distinct $V^+/V^-$ signatures (solar: 0.29, anthropogenic: 0.11, volcanic: 0.01) is a novel empirical finding.
It suggests that covariance-based fingerprinting can provide mechanistic information beyond what mean-field methods offer.
For example, two forcings that produce identical mean-field responses could in principle have distinguishable covariance fingerprints with different $V^+/V^-$ ratios.

\paragraph{Limitation.}
The $V^+/V^-$ decomposition is basis-dependent.
In the teleconnection analysis ($d=8$), the basis consists of physically meaningful climate indices, and the trace genuinely reflects total variance change.
In the gridded analysis ($d=58$), PCA dimensionality reduction precedes regression, and the projected regression coefficients live in a rotated basis with no simple physical interpretation.
The ``trace'' of the projected $\bm{\beta}$ is not the energy budget.
We recommend interpreting $V^+/V^-$ only when the covariance basis has clear physical meaning, and we restrict our mechanistic claims to the teleconnection-based results.

\subsection{Why attribution fractions are inherently unstable}\label{sec:attrib_discuss}

Attribution fractions---the percentage of explained covariance variance assigned to each forcing---are a natural quantity to report but are fundamentally less robust than detection $p$-values or $V^+/V^-$ ratios.
We elaborate on the three reasons given in Section~\ref{sec:attribution}.

\paragraph{Forcing interdependence.}
Anthropogenic warming changes the background state on which volcanic and solar forcings act.
A warmer ocean has different ENSO dynamics, so a volcanic eruption's teleconnection disruption depends on the anthropogenic trend.
The regression assumes additive forcing contributions, but the real system has interactions.
Varying the order of forcing entry (stepwise regression) changes fractions by 5--10 percentage points.

\paragraph{Non-additivity on the manifold.}
Covariance perturbations compose multiplicatively: if forcing~A sends $C$ to $A C A^\top$ and forcing~B sends $C$ to $B C B^\top$, the combined effect is $B A C A^\top B^\top \neq A C A^\top + B C B^\top - C$.
The tangent-space linearization makes effects approximately additive near the Fr\'{e}chet mean, but this is an approximation.
Attribution fractions inherit the linearization error, while the detection test (is the signal nonzero?) does not.

\paragraph{Methodological dependence.}
The fractions depend on window size (48 vs.\ 60 vs.\ 72 months shifts fractions by approximately 10 percentage points), shrinkage parameter ($\lambda_\mathrm{LW} = 0.10$ vs.\ $0.25$), regression method, and forcing normalization.
The sensitivity analyses in Appendix~\ref{app:sensitivity} document this instability.
The asymmetry is stark: detection $p$-values are robust across all tested methodological choices (all forcings remain significant at $\pval < 0.05$ for window sizes 48--84 months and shrinkage 0.05--0.30), while fractions are not.
This asymmetry is the core reason we frame the paper around detection and mechanism, not attribution percentages.

\subsection{Other limitations}\label{sec:limitations}

\paragraph{Low $R^2$.}
The model explains approximately 27--30\% of tangent-space variance.
This is not a failure of the method---internal variability (ENSO, weather noise, decadal oscillations) dominates climate covariance on interannual timescales.
The relevant point is that the 27--30\% that \emph{is} explained is attributable to specific forcings with high statistical confidence.
Stated differently, $R^2 \approx 0.30$ means 70\% of covariance variability is internal, which is entirely expected for a 5-year rolling window analysis dominated by ENSO cycling.

\paragraph{Rolling window assumptions.}
The 60-month window is a compromise: shorter windows are noisier (ill-conditioned covariance estimates), while longer windows smooth out volcanic signals.
Results are qualitatively robust to window sizes of 48--72 months but quantitatively sensitive (Table~\ref{tab:B1}).

\paragraph{Shrinkage.}
Ledoit--Wolf shrinkage toward a scaled identity biases covariance estimates toward isotropy.
This is necessary for positive-definiteness but may attenuate the $V^-$ signal.
The shrinkage parameter $\lambda_\mathrm{LW} = 0.15$ was not tuned; results are stable for $\lambda_\mathrm{LW} \in [0.10, 0.25]$ (Table~\ref{tab:B2}).

\paragraph{Solar significance.}
The 1983--2024 analysis window contains only approximately 4 solar cycles.
While solar forcing is detected at $\pval = 0.010$, this result should be treated with more caution than the anthropogenic and volcanic detections.
Extended windows (Table~\ref{tab:B3}) confirm solar significance with up to 7 cycles ($\pval = 0.014$), but the small number of cycles remains a limitation.
Solar detection would benefit most from extension to longer observational records.

\paragraph{No climate model comparison.}
Standard D\&A uses climate model ensembles (e.g., CMIP6 historical simulations) to define expected response patterns and estimate internal variability.
Our model-free approach eliminates model dependence but also foregoes the physical constraints that model-predicted patterns provide.
Validation against CMIP6 single-forcing experiments (historicalNat, historicalGHG, historicalVol) would provide the strongest possible test of Riemannian fingerprinting: in controlled model experiments, the true forcing is known, and one can verify that $V^+/V^-$ predictions match the actual forcing mechanism.
This is the most important direction for future work.

\subsection{Comparison with existing literature}\label{sec:literature}

Standard optimal fingerprinting \citep{Hasselmann1979,AllenStott2003,Hegerl2007} operates on spatial/temporal mean fields and is complementary to our covariance-based approach---the two methods detect different aspects of the forced climate response.
\citet{Ribes2013,Ribes2017} address the high-dimensionality challenge in standard fingerprinting via regularization, but still operate on means.

SPD matrices have been used in paleoclimate reconstructions \citep{Smerdon2015,Guillot2015}, but not for D\&A.
Tangent-space methods on SPD manifolds are well-established in brain--computer interfaces \citep{Barachant2013}, where they classify covariance matrices of EEG signals---the same mathematical framework applied to a different domain.
The Riemannian statistical foundations we build on are due to \citet{Pennec2006} and \citet{Said2017}.

To our knowledge, this is the first application of Riemannian tangent-space regression for climate forcing detection and mechanistic separation.

\subsection{Connection to the DeWitt supermetric}\label{sec:dewitt}

The affine-invariant metric~\eqref{eq:metric} on $\SPD(d)$ is a special case of the DeWitt supermetric on the space of Riemannian metrics $\GL^+(d)/\SO(d)$ \citep{DeWitt1967}.
The DeWitt metric is parametrized by a single constant $\lambda$:
\begin{equation}\label{eq:dewitt}
G_{ijkl}^{(\lambda)} = g_{ik}\, g_{jl} + g_{il}\, g_{jk} + \lambda\, g_{ij}\, g_{kl}.
\end{equation}
At $\lambda = -2/d$, the $V^+$ (conformal) mode decouples from $V^-$ (volume-preserving deformations).
Our $V^+/V^-$ decomposition corresponds exactly to this conformal/non-conformal split in the tangent space of $\SPD(d)$.

This connection is conceptually interesting but not essential to the results.
We note it because it places climate covariance analysis within a broader geometric framework that includes applications in general relativity \citep{DeWitt1967}, materials science (elastic tensor manifolds), and neuroscience (diffusion tensor imaging).
The same Riemannian machinery---Fr\'{e}chet mean, log map, tangent-space regression, $V^+/V^-$ decomposition---applies across these domains, suggesting that the approach may be of independent interest beyond climate science.

\subsection{Future directions}\label{sec:future}

Several extensions of this work are natural.

\paragraph{CMIP6 ground truth validation.}
Applying Riemannian fingerprinting to CMIP6 historical simulations (e.g., CESM2, UKESM, GFDL-ESM4), where the true forcing is known, would validate $V^+/V^-$ predictions against controlled experiments.
Single-forcing experiments (historicalNat, historicalGHG) could serve as the strongest test of mechanistic separation.
This requires approximately 1--10~GB per model from the Earth System Grid Federation.

\paragraph{Longer records.}
Extension back to 1900 using HadCRUT5, accepting sparser high-latitude coverage, would provide approximately 11 solar cycles for more robust solar detection and a longer anthropogenic trend for improved signal-to-noise.

\paragraph{Model-informed Riemannian fingerprinting.}
Using CMIP6 single-forcing experiments to define expected covariance response patterns as priors on the SPD manifold would combine the geometric advantages of our method with the physical constraints from climate models.

\paragraph{Spatio-temporal covariance.}
Extension to tensor product $\SPD(d_\mathrm{spatial}) \times \SPD(d_\mathrm{temporal})$ would detect forcings that change both spatial correlation patterns and temporal persistence structures.

\paragraph{Tipping point detection.}
Sudden acceleration along geodesics on $\SPD(d)$ could serve as an early warning indicator for climate tipping points (AMOC collapse, ice sheet instability), with $V^+/V^-$ identifying the underlying mechanism.


%% ====================================================================
\section{Conclusions}\label{sec:conclusions}
%% ====================================================================

We make three claims, supported by the results above.

\paragraph{1.\ Detection with geometric advantage.}
Riemannian fingerprinting detects all three major climate forcings (anthropogenic, volcanic, solar) as statistically significant ($\pval < 0.01$) in observed climate covariance data.
Standard Euclidean methods on the same data, with the same surrogates and forcing series, fail to reach significance for any forcing.
The geometric advantage arises from respecting the curved SPD manifold structure, exploiting the higher-dimensional signal space of second-order statistics (36 covariance DOF vs.\ 8 mean DOF for $d=8$), and linearizing at the intrinsic Fr\'{e}chet mean.
We note that this advantage is partly due to the increased dimensionality and partly due to the Riemannian geometry itself (Table~\ref{tab:threeway}).

\paragraph{2.\ Mechanistic separation via $V^+/V^-$.}
The $V^+/V^-$ decomposition---splitting regression coefficients into trace (energy budget) and traceless (teleconnection reorganization) components---reveals that the three forcings act through qualitatively distinct mechanisms.
Solar forcing operates primarily through the energy budget ($V^+ = 0.29$), anthropogenic forcing through a combination dominated by teleconnection reshaping ($V^+ = 0.11$), and volcanic forcing almost entirely through teleconnection reorganization ($V^+ = 0.01$).
This mechanistic fingerprint is a novel diagnostic unavailable in standard mean-field D\&A.

\paragraph{3.\ Model-free approach.}
The method requires no climate model ensembles---only observed data and forcing time series.
This eliminates model dependence (a persistent concern in standard D\&A) at the cost of not having model-predicted response patterns as priors.

\paragraph{Caveats.}
Attribution fractions can be computed but are inherently imprecise: forcings are interdependent, effects are non-additive on the covariance manifold, and fractions depend on methodological choices that do not affect detection significance.
Explained variance is approximately 30\%, consistent with dominant internal variability.
The method detects forcings in covariance structure, which is complementary to---not a replacement for---standard mean-field D\&A.
Validation against CMIP6 controlled experiments is the most important next step.

The mathematical framework (tangent-space regression on SPD manifolds) is general and applicable to any domain where covariance matrices are the objects of interest, including neuroscience (EEG/fMRI connectivity), finance (portfolio risk), and materials science (elastic tensor fields).


%% ====================================================================
%% APPENDICES
%% ====================================================================
\appendix

\makeatletter
\renewcommand{\thefigure}{\thesection\arabic{figure}}
\renewcommand{\thetable}{\thesection\arabic{table}}
\@addtoreset{figure}{section}
\@addtoreset{table}{section}
\makeatother

\section{Mathematical Details}\label{app:math}

\subsection{Norm preservation under vech vectorization}

The $\vech$ operator extracts the lower-triangular entries of a symmetric matrix.
To preserve the Frobenius inner product under vectorization, off-diagonal elements must be scaled by $\sqrt{2}$.
Specifically, for $V, W \in \Sym(d)$,
\begin{equation}
\tr(V W) = \sum_{i} V_{ii} W_{ii} + 2 \sum_{i < j} V_{ij} W_{ij} = \vech(V)^\top \vech(W),
\end{equation}
where $\vech$ maps $V_{ii} \mapsto V_{ii}$ and $V_{ij} \mapsto \sqrt{2}\, V_{ij}$ for $i \neq j$.
This ensures that the multivariate regression in $\RR^{d(d+1)/2}$ respects the Riemannian inner product in the tangent space.

\subsection{Fr\'{e}chet mean convergence}

On $\SPD(d)$ with the affine-invariant metric, the Fr\'{e}chet mean exists and is unique for any finite collection of SPD matrices \citep{Moakher2005}.
The Riemannian gradient descent algorithm converges at a linear rate with step size $\tau = 1$ when the data lie within a geodesically convex ball.
In practice, the log-Euclidean initialization $\exp\!\bigl(\frac{1}{N}\sum_i \log C_i\bigr)$ provides a warm start that lies close to the true Fr\'{e}chet mean, and convergence to $\|\nabla\|_F < 10^{-8}$ is achieved within 5--10 iterations for all datasets in this study.

\subsection{Block-permutation vs.\ phase-randomization surrogates}

Block-permutation surrogates preserve the temporal structure \emph{within} blocks of length $b$ while destroying the forcing--response coupling.
Phase-randomization surrogates preserve the full power spectrum by randomizing Fourier phases.
Both methods yield similar $p$-values in our analysis (all within a factor of 2).
We prefer block permutation because (i)~it makes fewer assumptions about the spectral structure, (ii)~it is straightforward to implement and interpret, and (iii)~the block size provides an explicit knob for sensitivity testing (Appendix~\ref{app:blocksize}).

\section{Sensitivity Analysis}\label{app:sensitivity}

All sensitivity analyses use 8 teleconnection indices + GMTA, GloSSAC v2.22 volcanic AOD, and the 1983--2024 window unless otherwise noted.

\begin{figure}[h]
\centering
\includegraphics[width=\textwidth]{figures/fig_sensitivity.pdf}
\caption{Sensitivity of detection significance ($-\log_{10}(p)$) to window size (left), shrinkage parameter (center), and surrogate block size (right). Dashed lines mark $p=0.05$; dotted lines mark $p=0.01$. All three forcings remain above the $p=0.05$ threshold for window sizes 48--84 months and shrinkage $\lambda \in [0.05, 0.30]$. Significance degrades at block sizes $\geq 12$ due to insufficient permutation resolution (fewer than 13 permutable blocks).}
\label{fig:sensitivity}
\end{figure}

\subsection{Sensitivity to rolling window size}\label{app:window}

\begin{table}[h]
\centering
\caption{Sensitivity of detection $p$-values and $V^+$ ratio to rolling window size. All three forcings remain significant ($\pval < 0.05$) at every tested window size.}
\label{tab:B1}
\begin{tabular}{lccccc}
\toprule
Window (months) & $\pval_\mathrm{anthro}$ & $\pval_\mathrm{volc}$ & $\pval_\mathrm{solar}$ & $R^2$ & $V^+_\mathrm{solar}$ \\
\midrule
48 & $0.002^{**}$ & $0.002^{**}$ & $0.008^{**}$ & 0.231 & 0.332 \\
60 & $0.002^{**}$ & $0.002^{**}$ & $0.010^{**}$ & 0.276 & 0.294 \\
72 & $0.004^{**}$ & $0.002^{**}$ & $0.006^{**}$ & 0.286 & 0.328 \\
84 & $0.008^{**}$ & $0.002^{**}$ & $0.014^{*}$ & 0.299 & 0.380 \\
\bottomrule
\multicolumn{6}{l}{\footnotesize $^{**}\pval < 0.01$;\; $^{*}\pval < 0.05$.}
\end{tabular}
\end{table}

All three forcings remain significant ($\pval < 0.05$) at every window size from 48 to 84 months.
The $V^+$ ratio for solar forcing is stable at 0.29--0.38, confirming the energy-budget mechanism across window choices.
Longer windows produce marginally higher $R^2$ (more temporal smoothing) but slightly weaker detection of the anthropogenic trend, likely due to reduced temporal resolution.

\subsection{Sensitivity to Ledoit--Wolf shrinkage}\label{app:shrinkage}

\begin{table}[h]
\centering
\caption{Sensitivity to Ledoit--Wolf shrinkage parameter $\lambda_\mathrm{LW}$. All forcings remain significant across the full range. Higher shrinkage biases toward isotropy, slightly weakening detection.}
\label{tab:B2}
\begin{tabular}{lcccc}
\toprule
$\lambda_\mathrm{LW}$ & $\pval_\mathrm{anthro}$ & $\pval_\mathrm{volc}$ & $\pval_\mathrm{solar}$ & $R^2$ \\
\midrule
0.05 & $0.002^{**}$ & $0.002^{**}$ & $0.006^{**}$ & 0.294 \\
0.10 & $0.002^{**}$ & $0.002^{**}$ & $0.006^{**}$ & 0.283 \\
0.15 & $0.002^{**}$ & $0.002^{**}$ & $0.010^{**}$ & 0.276 \\
0.20 & $0.002^{**}$ & $0.002^{**}$ & $0.012^{*}$ & 0.269 \\
0.30 & $0.004^{**}$ & $0.002^{**}$ & $0.012^{*}$ & 0.257 \\
\bottomrule
\multicolumn{5}{l}{\footnotesize $^{**}\pval < 0.01$;\; $^{*}\pval < 0.05$.}
\end{tabular}
\end{table}

Higher shrinkage biases covariance estimates toward isotropy (scaled identity), attenuating the $V^-$ signal and slightly weakening detection.
Nevertheless, all three forcings remain significant at $\pval < 0.05$ across the full range $\lambda_\mathrm{LW} \in [0.05, 0.30]$.

\subsection{Extended time window (solar robustness)}\label{app:extended}

\begin{table}[h]
\centering
\caption{Solar detection robustness with extended time windows. With more solar cycles, solar remains significant but $V^+$ decreases as teleconnection effects ($V^-$) accumulate over longer periods.}
\label{tab:B3}
\begin{tabular}{lcccccc}
\toprule
Period & Months & Solar cycles & $\pval_\mathrm{solar}$ & $V^+_\mathrm{solar}$ & $\pval_\mathrm{anthro}$ & $\pval_\mathrm{volc}$ \\
\midrule
1950--2024 & 804 & $\sim$7 & $0.014^{*}$ & 0.215 & $0.002^{**}$ & $0.002^{**}$ \\
1958--2024 & 804 & $\sim$6 & $0.014^{*}$ & 0.215 & $0.002^{**}$ & $0.002^{**}$ \\
1970--2024 & 660 & $\sim$5 & $0.032^{*}$ & 0.264 & $0.002^{**}$ & $0.002^{**}$ \\
1983--2024 & 504 & $\sim$4 & $0.010^{**}$ & 0.294 & $0.002^{**}$ & $0.002^{**}$ \\
\bottomrule
\multicolumn{7}{l}{\footnotesize $^{**}\pval < 0.01$;\; $^{*}\pval < 0.05$.}
\end{tabular}
\end{table}

Solar forcing remains significant ($\pval < 0.05$) with up to 7 solar cycles.
However, only 4 cycles fall within the primary analysis window (1983--2024), and the solar $p$-values do not uniformly improve with longer windows---reflecting the tension between more cycles and noisier/sparser early data.
The $V^+$ ratio for solar decreases with longer windows (0.29 to 0.22), suggesting that teleconnection effects ($V^-$) accumulate over longer periods relative to the energy-budget signal.

\subsection{Three-way method comparison}\label{app:threeway}

Table~\ref{tab:threeway} (reproduced in the main text) isolates the geometric effect from the dimensionality effect.
The flat Euclidean control (vectorized covariance matrices without Riemannian geometry) improves over mean-field Euclidean regression for volcanic and anthropogenic forcing simply due to the larger signal space (36D vs.\ 8D).
The additional improvement from flat to Riemannian---at the \emph{same} 36 dimensions---demonstrates an independent geometric contribution: the log map and Fr\'{e}chet mean provide a better linearization of the covariance dynamics than naive vectorization.

This decomposition is important for intellectual honesty: the Riemannian advantage is \emph{not} purely geometric.
Part of it comes from working with second-order statistics (more dimensions), and part comes from the correct geometry on those statistics.
Both contributions are real and complementary.

\subsection{Sensitivity to surrogate block size}\label{app:blocksize}

\begin{table}[h]
\centering
\caption{Sensitivity of detection $p$-values to surrogate block size (in years). Block size controls the trade-off between preserving autocorrelation (larger blocks) and permutation resolution (more blocks).}
\label{tab:B5}
\begin{tabular}{lccccc}
\toprule
Block size (yr) & $\pval_\mathrm{anthro}$ & $\pval_\mathrm{volc}$ & $\pval_\mathrm{solar}$ & Permutable blocks \\
\midrule
2 & $0.002^{**}$ & $0.002^{**}$ & $0.002^{**}$ & 75 \\
3 & $0.002^{**}$ & $0.002^{**}$ & $0.002^{**}$ & 50 \\
5 & $0.004^{**}$ & $0.002^{**}$ & $0.008^{**}$ & 30 \\
8 & $0.014^{*}$ & $0.002^{**}$ & $0.028^{*}$ & 19 \\
12 & $0.066$ & $0.004^{**}$ & $0.082$ & 13 \\
20 & $0.270$ & $0.026^{*}$ & $0.275$ & 8 \\
\bottomrule
\multicolumn{5}{l}{\footnotesize $^{**}\pval < 0.01$;\; $^{*}\pval < 0.05$.}
\end{tabular}
\end{table}

Block size controls the trade-off between preserving autocorrelation within blocks (larger blocks are more conservative) and having enough permutable units for a well-resolved null distribution (more blocks provide finer $p$-value resolution).

At block sizes $\geq 12$~years, anthropogenic and solar lose significance.
This reflects insufficient permutation resolution (only 13 blocks yield a coarse null distribution), not absence of signal.
The critical point is that significance depends on block size choice, and we are transparent about this dependence.
At block sizes of 5--8 years, all three forcings are significant ($\pval < 0.05$) while preserving substantial within-block autocorrelation.

Volcanic forcing is the most robust: it remains significant at $\pval < 0.05$ at all block sizes up to 20 years, likely because the episodic volcanic signal (Pinatubo, El~Chich\'{o}n) is so distinctive that even a coarse null distribution cannot produce comparable spurious signals.

We report block size 5~years as the primary result because it provides 30 permutable blocks---sufficient for a well-resolved null---while preserving autocorrelation on the timescale of the rolling window.


%% ====================================================================
%% AVAILABILITY STATEMENT
%% ====================================================================

\section*{Availability Statement}

All observational data used in this study are publicly available.
Teleconnection indices were obtained from the NOAA Climate Prediction Center (https://www.cpc.ncep.noaa.gov/).
HadCRUT5 gridded temperatures were obtained from the Met Office Hadley Centre (https://www.metoffice.gov.uk/hadobs/hadcrut5/).
CO$_2$ and CH$_4$ concentration data were obtained from the NOAA Global Monitoring Laboratory (https://gml.noaa.gov/).
GloSSAC v2.22 stratospheric aerosol optical depth data were obtained from NASA Langley (https://doi.org/10.5067/GloSSAC-L3-V2.22).
Sunspot number data were obtained from SILSO (https://www.sidc.be/silso/).
The implementation uses \texttt{numpy} and \texttt{scipy}; no climate-specific libraries are required.
Code is available at https://github.com/sausterm/omnisciences-research.


%% ====================================================================
%% REFERENCES
%% ====================================================================

\begin{thebibliography}{99}

\bibitem[Adams et~al.(2003)]{Adams2003}
Adams, J.~B., Mann, M.~E., \& Ammann, C.~M. (2003).
Proxy evidence for an El Ni\~{n}o-like response to volcanic forcing.
\textit{Nature}, \textit{426}, 274--278.

\bibitem[Allen and Stott(2003)]{AllenStott2003}
Allen, M.~R. \& Stott, P.~A. (2003).
Estimating signal amplitudes in optimal fingerprinting, Part~I: Theory.
\textit{Climate Dynamics}, \textit{21}, 477--491.

\bibitem[Barachant et~al.(2013)]{Barachant2013}
Barachant, A., Bonnet, S., Congedo, M., \& Jutten, C. (2013).
Classification of covariance matrices using a Riemannian-based kernel for BCI applications.
\textit{Neurocomputing}, \textit{112}, 172--178.

\bibitem[Cohen et~al.(2014)]{Cohen2014}
Cohen, J., Screen, J.~A., Furtado, J.~C., Barlow, M., Whittleston, D., Coumou, D., Francis, J., Dethloff, K., Entekhabi, D., Overland, J., \& Jones, J. (2014).
Recent Arctic amplification and extreme mid-latitude weather.
\textit{Nature Geoscience}, \textit{7}, 627--637.

\bibitem[DeWitt(1967)]{DeWitt1967}
DeWitt, B.~S. (1967).
Quantum theory of gravity.~I. The canonical theory.
\textit{Physical Review}, \textit{160}(5), 1113--1148.

\bibitem[Francis and Vavrus(2012)]{FrancisVavrus2012}
Francis, J.~A. \& Vavrus, S.~J. (2012).
Evidence linking Arctic amplification to extreme weather in mid-latitudes.
\textit{Geophysical Research Letters}, \textit{39}, L06801.

\bibitem[Gray et~al.(2010)]{Gray2010}
Gray, L.~J., Beer, J., Geller, M., Haigh, J.~D., Lockwood, M., Matthes, K., Cubasch, U., Fleitmann, D., Harrison, G., Hood, L., Luterbacher, J., Meehl, G.~A., Shindell, D., van Geel, B., \& White, W. (2010).
Solar influences on climate.
\textit{Reviews of Geophysics}, \textit{48}, RG4001.

\bibitem[Guillot et~al.(2015)]{Guillot2015}
Guillot, D., Raber, B., \& Sun, Q. (2015).
Statistical analysis of positive definite matrices with applications in diffusion tensor imaging.
\textit{Annals of Applied Statistics}, \textit{9}(2), 901--923.

\bibitem[Hasselmann(1979)]{Hasselmann1979}
Hasselmann, K. (1979).
On the signal-to-noise problem in atmospheric response studies.
In B.~D.~Shaw (Ed.), \textit{Meteorology of Tropical Oceans} (pp.\ 251--259).
Royal Meteorological Society.

\bibitem[Hegerl et~al.(2007)]{Hegerl2007}
Hegerl, G.~C., Zwiers, F.~W., Braconnot, P., Gillett, N.~P., Luo, Y., Marengo Orsini, J.~A., Nicholls, N., Penner, J.~E., \& Stott, P.~A. (2007).
Understanding and attributing climate change.
In \textit{Climate Change 2007: The Physical Science Basis. Contribution of Working Group~I to the Fourth Assessment Report of the IPCC} (Ch.~9).

\bibitem[{IPCC}(2021)]{IPCC_AR6_Ch3}
IPCC (2021).
Human influence on the climate system.
In V.~Masson-Delmotte et~al.\ (Eds.), \textit{Climate Change 2021: The Physical Science Basis. Contribution of Working Group~I to the Sixth Assessment Report of the IPCC} (Ch.~3).
Cambridge University Press.

\bibitem[Moakher(2005)]{Moakher2005}
Moakher, M. (2005).
A differential geometric approach to the geometric mean of symmetric positive-definite matrices.
\textit{SIAM Journal on Matrix Analysis and Applications}, \textit{26}(3), 735--747.

\bibitem[Pennec et~al.(2006)]{Pennec2006}
Pennec, X., Fillard, P., \& Ayache, N. (2006).
A Riemannian framework for tensor computing.
\textit{International Journal of Computer Vision}, \textit{66}(1), 41--66.

\bibitem[Predybaylo et~al.(2017)]{Predybaylo2017}
Predybaylo, E., Stenchikov, G.~L., Wittenberg, A.~T., \& Zeng, F. (2017).
El Ni\~{n}o/Southern Oscillation response to low-latitude volcanic eruptions depends on ocean pre-conditions and eruption timing.
\textit{Communications Earth \& Environment}.

\bibitem[Ribes et~al.(2013)]{Ribes2013}
Ribes, A., Planton, S., \& Terray, L. (2013).
Application of regularised optimal fingerprinting to attribution.
Part~I: Method, properties, and idealised analysis.
\textit{Climate Dynamics}, \textit{41}, 2817--2836.

\bibitem[Ribes et~al.(2017)]{Ribes2017}
Ribes, A., Zwiers, F.~W., Azais, J.-M., \& Naveau, P. (2017).
A new statistical approach to climate change detection and attribution.
\textit{Climate Dynamics}, \textit{48}, 367--386.

\bibitem[Said et~al.(2017)]{Said2017}
Said, S., Bombrun, L., Berthoumieu, Y., \& Manton, J.~H. (2017).
Riemannian Gaussian distributions on the space of symmetric positive definite matrices.
\textit{IEEE Transactions on Information Theory}, \textit{63}(4), 2153--2170.

\bibitem[Scaife et~al.(2013)]{Scaife2013}
Scaife, A.~A., Ineson, S., Knight, J.~R., Gray, L., Kodera, K., \& Smith, D.~M. (2013).
A mechanism for lagged North Atlantic climate response to solar variability.
\textit{Geophysical Research Letters}, \textit{40}, 434--439.

\bibitem[Smerdon et~al.(2015)]{Smerdon2015}
Smerdon, J.~E., Coats, S., \& Ault, T.~R. (2015).
Model-dependent spatial skill in pseudoproxy experiments testing climate field reconstruction methods for the Common Era.
\textit{Climate Dynamics}, \textit{46}, 1921--1942.

\end{thebibliography}

\end{document}
